% DOC TYPE E DEF QB SOFTWARE %%%%%%%%%%%%%%%%%%%%%%%%%
\documentclass[12pt]{article}
%%%%%%%%%%%%%%%%%%%% PACKAGES %%%%%%%%%%%%%%%%%%%%
\usepackage[english,italian]{babel}
\usepackage[a4paper, margin=3cm]{geometry}
\usepackage[T1]{fontenc}
\usepackage[utf8]{inputenc}
\usepackage[table]{xcolor}
\usepackage{amsmath}
\usepackage{graphicx}
\usepackage{float}
\usepackage{tabularx}
\usepackage{booktabs}
\usepackage{hyperref}
\usepackage{xcolor}
\usepackage{multicol}
\usepackage{multirow}
\usepackage{soul}
\usepackage{enumitem}
\usepackage{textcomp}
\usepackage{eurosym}
\usepackage{lastpage}
\usepackage{fancyhdr}
\usepackage{adjustbox}
\usepackage{subfiles}

%%%%%%%%%%%%%%%%%%%% COMMAND -> New commands %%%%%%%%%%%%%%%%%%%%
\newcommand{\mailtoQBS}
{
	\href{mailto:qbsoftware.swe@gmail.com}{qbsoftware.swe@gmail.com}
}

\let\oldpar\paragraph
\renewcommand{\paragraph}[1]{\oldpar{#1}\mbox{}\\}

%%%%%%%%%%%%%%%%%%%% COMMAND -> Redefine LaTeX Commands %%%%%%%%%%%%%%%%%%%%
\def\title#1{\gdef\THETITLE{#1}}
\def\date#1{\gdef\THEDATE{#1}}
\def\footname#1{\gdef\THEFOOTNAME{#1}}

\let\oldtexteuro\texteuro
\renewcommand{\texteuro}{\euro}

%%%%%%%%%%%%%%%%%%%% STYLES -> Commands %%%%%%%%%%%%%%%%%%%%
\newcommand{\makefirstpage}
{
	\begin{titlepage}
		
		% Defines a new command for the horizontal lines, change thickness here
		\newcommand{\HRule}{\rule{\linewidth}{0.2mm}} 
		
		\center % Center everything on the page
		
		
		%	Heading Sections
		
		\textsc{\LARGE QB Software}\\[.1cm] 
		\includegraphics[scale=.15]{imgs/qb-software-logo.png}\\[-.1 cm] 
		\includegraphics[scale=.025]{imgs/x.png}\\[0.5cm]
		\includegraphics[scale=.3]{imgs/unipd_logo.png}\\[.5cm]
		\textsc{\Large Università degli studi di Padova}\\[0.5cm] 
		\textsc{\large corso di ingegneria del software }\\[0.5cm]
		\textsc{\large anno accademico 2023/2024 }\\[0.5cm]
		
		
		%	Title section and date section
		
		\ifdefined\THEDATE
		\HRule \\[0.4cm]
		{ \huge{ \bfseries {\THETITLE}} \\ [.5cm]
			\THEDATE}\\[0.4cm] 
		\HRule \\[0.4cm]
		\else
		\HRule \\[0.4cm]
		{ \huge{ \bfseries {\THETITLE}} } \\
		\HRule \\[0.4cm]
		\fi
		%
		\textsc{Contatti:} \mailtoQBS\\[0.3cm]
		
		\vfill 
		
	\end{titlepage}
}

%%%%%%%%%%%%%%%%%%%% ACCESSIBILITY %%%%%%%%%%%%%%%%%%%%
\let\oldhref\href
\renewcommand{\href}[2]{\oldhref{#1}{\textcolor{blue}{\ul{#2}}}}

\hypersetup{
	colorlinks = true,
	linkcolor = cyan,
}

%%%%%%%%%%%%%%%%%%%% STYLES -> Settings %%%%%%%%%%%%%%%%%%%%
% Enable header and footer style
\pagestyle{fancy}
\thispagestyle{empty}
\thispagestyle{fancy}
\pagestyle{fancy}

% Header
\setlength{\topmargin}{-40pt}
\setlength{\headsep}{60pt}
\fancyhf{}
\lhead{QB Software}
\setlength{\headheight}{15pt}
\rhead{\includegraphics[width=1cm]{imgs/qb-software-logo.png}}

% Footer 
\fancyfoot{}
\fancyfoot[L]{\THEFOOTNAME}
\fancyfoot[R]{Pagina \thepage~di~\pageref{LastPage}}
\futurelet\TMPfootrule\def\footrule{\TMPfootrule}
\setcounter{page}{0}
\pagenumbering{arabic}
\renewcommand{\footrulewidth}{0.3pt}

%%%%%%%%%%%%%%%%%%%% ENVIRONMENT %%%%%%%%%%%%%%%%%%%%
% Remove colorless padding in booktabs
\setlength{\aboverulesep}{0cm}
\setlength{\belowrulesep}{0cm}
\setlength{\extrarowheight}{.75ex}

\newenvironment{todo}{
	\rowcolors{2}{cyan!80!black!30!}{cyan!80!black!20!}
	\begin{tabular}{p{3.48cm}>{\raggedright\arraybackslash}p{4cm}>{\raggedright\arraybackslash}p{6.5cm}}
		\toprule
		\rowcolor{gray!20} \textbf{ID}	& \textbf{Interessato} & \textbf{Task} 
		\\\midrule
	}{
		\bottomrule
	\end{tabular}
}

\newenvironment{changelog}{
	\noindent
	{\Large \textbf{Registro delle modifiche}}
	\noindent
	\begin{table}[h]
		\rowcolors{2}{cyan!80!black!30!}{cyan!80!black!20!}
		\begin{adjustbox}{width=\textwidth}
			\begin{tabular}{|c|c|p{2.35cm}|c|p{3.3cm}|}
				\hline
				\rowcolor{gray!20}
				\textbf{V.} & \textbf{Data} & \textbf{Membro} & \textbf{Ruolo} & \textbf{Descrizione} \\
				
				\hline
			}{
				\hline
			\end{tabular}
		\end{adjustbox}
	\end{table}
	
	\clearpage
}

% 1   2    3        4             5            6            7          
% Ver Data Relatore RuoloRelatore DataVerifica Verificatore Descrizione
\newcommand{\newlog}[7]{
	#1 & #5 & #6 & Verificatore & Controllo qualità \\
	   & #2 & #3 & #4           & #7 \\\hline
}

% INFORMAZIONI DOCUMENTO %%%%%%%%%%%%%%%%%%%%%%%%%
\title{Verbale interno}
\date{del 18/12/2023}
\footname{Verbale interno del 18/12/2023}

\begin{document}
	% PRIME PAGINE %%%%%%%%%%%%%%%%%%%%%%%%%
	\makefirstpage
	
	% 1   2    3        4             5            6            7          
% Ver Data Relatore RuoloRelatore DataVerifica Verificatore Descrizione
\begin{changelog}
	\newlog{0.3.0}{15/12/2023}{S. Rovea}{Amministratore}{/12/2023}{}{Correzione ed espansione sezione \ref{sec:management_process}}
	\newlog{0.2.0}{02/12/2023}{S. Destro}{Amministratore}{03/12/2023}{A. Feltrin}{Aggiunta sezione su processi organizzativi, sezione \ref{sec:management_process}}
	\newlog{0.1.0}{22/11/2023}{A. Bustreo}{Amministratore}{22/11/2023}{A. Giurisato}{Aggiunto processo supporto per la documentazione, sezione \ref{sec:documentation}. Aggiunta prefazione, sezione \ref{sec:prefazione}}
\end{changelog}
	\clearpage
	
	\tableofcontents
	\clearpage

    \section{Informazioni generali}
	
	\subsection{Luogo e data dell'incontro}
	
    	\begin{itemize}
    		\item \textbf{Luogo}: meeting su Discord
    		\item \textbf{Data}: 18/12/2023
    		\item \textbf{Ora di inizio}: 16:00
    		\item \textbf{Ora di fine}: 18:00
    	\end{itemize}
	
	\subsection{Presenze}
	
    	\begin{itemize}
    		\item \textbf{Totale presenze}:
    		\begin{itemize}
    			\item Bustreo Alessandro
    			\item Destro Stefano
    			\item Domuta Alessia 
    			\item Feltrin Alessandro 
    			\item Fontana Raffaele Paolo 
    			\item Giurisato Andrea 
    			\item Rovea Silvia
    		\end{itemize}
    		
    		\item \textbf{Totale assenze}: 0
    		
    		\item \textbf{Partecipanti esterni}:
    		\begin{itemize}
    			\item nessuno
    		\end{itemize}
    	\end{itemize}

    \section{Ordine del giorno}
        Questioni fissate nella riunione interna precedente:
        \begin{itemize}
			\item pianificare le singole attività per l'inizio del prossimo Sprint (Sprint-3).
    	\end{itemize}
    	%
    	Nuove questioni:
    	\begin{itemize}
    		\item retrospettiva dello Sprint passato (Sprint-2).
    	\end{itemize}
    
    \section{Verbale}

	La riunione si è aperta discutendo la retrospettiva relativa allo Sprint-2, riportata dettagliatamente nella sezione \hyperref[sec:retrospettiva]{Retrospettiva}, seguita dalla pianificazione del prossimo Sprint (Sprint-3), riportata dettagliatamente nella sezione \hyperref[sec:pianificazione]{Pianificazione}.
	
	\vspace{0.3cm}
	\noindent
	In seguito è stato deciso di richiedere un incontro con il professor Cardin durante questo Sprint per presentare il nostro lavoro e i nostri dubbi riguardanti l'Analisi dei Requisiti.


	\subsection{Retrospettiva} \label{sec:retrospettiva}

	Si riporta di seguito la retrospettiva relativa al secondo Sprint (dal 04/12/2023 al 17/12/2023).

	\subsubsection{Cosa ha funzionato}
	\begin{itemize}
		\item Pianificazione attività: gli obiettivi del secondo Sprint sono stati tutti raggiunti;
		\item sistema di review di GitHub: il sistema di review di GitHub si è rivelato pratico e ci ha permesso di comunicare meglio tra di noi durante le attività di verifica;
		\item sistema di automazione: il sistema di automazione, in seguito alle modifiche discusse durante la riunione del 12/04/2023, è stato migliorato con l'aggiunta della possibilità di controllare lo stato delle action prima di confermare una pull request. Questo si è rivelato molto utile e ci permette di evitare numerosi errori;
		\item quantità del carico di lavoro equa: la quantità di ore produttive prevista è stata distribuita equamente tra i membri del gruppo.
	\end{itemize}
	
	\subsubsection{Cosa non ha funzionato}
	\begin{itemize}
		\item Distribuzione temporale del carico di lavoro dei Verificatori: il carico di lavoro dei Verificatori si è concentrato in particolare alla fine dello Sprint. Per il prossimo Sprint è necessario distribuire meglio il lavoro, così da terminare prima alcune issue e non concentrare il lavoro dei Verificatori a fine Sprint.
	\end{itemize}


	\subsubsection{Rischi}
	I rischi previsti durante la pianificazione del secondo sprint erano i seguenti: 
	\begin{itemize}
		\item RR1 - Incomprensione requisiti;
		\item RT1 - Inesperienza con le tecnologie.
	\end{itemize}
	I rischi effettivamente incontrati sono i seguenti:
	\begin{itemize}
		\item RR1 - Incomprensione requisiti: il documento Analisi dei Requisiti è stato modificato più volte. Tuttavia, poichè gli errori sono stati gestiti in una fase ancora iniziale del documento, non ci sono stati ritardi;
		\item RO2 - Impegni personali: un membro del team non è stato presente per una settimana. Tuttavia, il piano di controllo è stato seguito correttamente e non ci sono state conseguenze;
		\item RO5 - Reperibilità esterna: a causa del licenziamento del project manager, il gruppo ha avuto difficoltà nel fissare un incontro con l'azienda. Tuttavia, è stato applicato il piano di contingenza previsto e non ci sono state conseguenze.
	\end{itemize}
	Il piano di mitigazione si è rivelato efficace per tutti i rischi incontrati, dunque non si rendono necessarie attività di modifica della gestione dei rischi.

	\subsubsection{Tabella riassuntiva}

	\begin{table}[H]
		\begin{tabularx}{\textwidth}{X|X}
			\hline
			\multicolumn{1}{|>{\centering\arraybackslash}X|}{\textbf{Keep doing}}
			&
			\multicolumn{1}{>{\centering\arraybackslash}X|}{\textbf{Improvements}}
			\\\hline\hline
			Obiettivi del secondo Sprint raggiunti
			&
			Distribuzione temporale del carico di lavoro dei Verificatori
			\\\arrayrulecolor{gray}\hline
			Divisione equa del carico di lavoro nel team
			&

			\\\arrayrulecolor{gray}\hline
			L'utilizzo degli strumenti per la review di GitHub
			&

			\\\arrayrulecolor{gray}\hline
			Sistema di automazione
			&

		\end{tabularx}
		\caption{retrospettiva del 18/12/2023.}
	\end{table}

	\subsection{Pianificazione} \label{sec:pianificazione}
	\subsubsection{Obiettivi}
	Gli obiettivi per lo Sprint 3 (dal 18/12/2023 al 31/12/2023) decisi dal gruppo sono i seguenti:
	\begin{itemize}
		\item proseguire con lo studio delle tecnologie da utilizzare per la realizzazione del PoC, in particolare MongoDB, RESTEasy e JMAP;
		\item redigere una prima versione del documento Piano di Qualifica;
		\item continuare il documento Analisi dei Requisiti.
	\end{itemize}
	\subsubsection{Rischi}
	I rischi previsti sono i seguenti:
	\begin{itemize}
		\item RT1 - Inesperienza con le tecnologie: prevista a causa dell'inesperienza con le tecnologie che il gruppo pianifica di studiare in questo Sprint;
		\item RO5 - Reperibilità esterna: prevista a causa delle festività invernali;
		\item RO2 - Impegni personali: prevista a causa delle festività invernali;
	\end{itemize}
	
    
    \section{Azioni da intraprendere}
    
        \begin{todo}
            \href{https://github.com/QB-Software-swe/docs/issues/69}{VI-2023-12-18\#69}
            &
            R. Fontana
            &
            Prima stesura del documento Piano di Qualifica
            \\\midrule
            \href{https://github.com/QB-Software-swe/docs/issues/70}{VI-2023-12-18\#70}
            &
            S. Rovea 
            &
            Continuazione stesura Casi d'Uso nel documento Analisi dei Requisiti
            \\
			\midrule
            \href{https://github.com/QB-Software-swe/docs/issues/71}{VI-2023-12-18\#71}
            &
            A. Giurisato
            &
            Stesura sezione Requisiti nel documento Analisi dei Requisiti
            \\\midrule
            \href{https://github.com/QB-Software-swe/docs/issues/72}{VI-2023-12-18\#72}
            &
            A. Domuta
            &
            Continuazione RESTEasy configuration per PoC
            \\\midrule
            \href{https://github.com/QB-Software-swe/docs/issues/73}{VI-2023-12-18\#73}
            &
            S. Destro
            &
            Configurazione JMAP lato server per PoC        
            \\
    	\end{todo}
    
    \section{Ordine del giorno per la prossima riunione interna}
        \begin{itemize}
        		\item pianificare le attività per l'inizio dello Sprint-4 (dal 01/01/2024 al 07/01/2024)
        		\item retrospettiva dello Sprint-3 (dal 18/12/2023 al 31/12/2023)
    	\end{itemize}
\end{document}